\futureoutlook
\label{sec:future}

The five-stage pipeline now runs end to end, and inside the aluminium-rich
design region it reproduces experimental elastic data to within a few percent.
The weakness is everything outside that region, where the refit
potential turns mechanically unstable (Section~\ref{sec:nnip}). The remaining
work falls into three areas: making the optimization better posed, widening the
chemistry it can cover, and tightening the reference data behind it.

\subsection{Upcoming Objectives}

\paragraph{A better-posed optimization.}
The refit currently moves all $84$ free cross-term parameters at once,
which is more freedom than the elastic targets can constrain. The plan is to
restrict the search to the single-element-anchored cross-terms that actually
matter, ranked by a Sobol sensitivity analysis, which cuts the active
parameter count to a tractable few. The Born stability conditions,
applied as a hard filter only at the end now, should instead be baked into the
training loss as a penalty, so the surrogate learns to avoid unstable regions
rather than being rejected after the fact. The sampling base would also grow,
through active learning, to at least $500$ composition medoids.

\paragraph{Beyond the aluminium domain.}
Several target chemistries need elements the current basis does not carry, so
the basis would be extended to out-of-basis species (C, V, N) that the ferrous
and titanium alloys depend on. The MD and DFT sampling would be pushed
deliberately into the HCP region (Ti, Mg, Zn) and the ferrous region
(Fe--Cr--Ni), which the present corpus barely touches. A symmetry-aware
elastic extraction that uses the full hexagonal stiffness tensor, rather than
the cubic strain protocol, would remove a known source of error on the HCP
metals.

\paragraph{Reference data and validation.}
For the cross-terms, special quasi-random structure references
\citep{lewis1995} computed at intermediate compositions would give the fit
something to anchor to between the pure-element endpoints, and the DFT set
would be enlarged to span the broadened element domain. The ASM/MatWeb
comparison of Section~\ref{sec:validation} becomes a standing external check
that every refit round is measured against, rather than a one-off at the end.

\subsection{Updated Semester Plan}

The Gantt chart of Figure~\ref{fig:gantt} shows the completed schedule. All
eight work packages, from the literature review through the closed-loop
inverse optimization and the external validation, finished within the
semester, leaving the writing and poster preparation for the final weeks.

\begin{figure}[htbp]
	\centering
	\begin{semestergantt}
		\workpackage{Literature Review}{1}{3}
		\workpackage{DFT Reference Calculations}{2}{5}
		\workpackage{MEAM Initialization Pipeline}{4}{7}
		\workpackage{NN Surrogate Training}{6}{9}
		\workpackage{Inverse Optimization}{8}{13}
		\workpackage{Validation \& Testing}{10}{14}
		\workpackage{Writing \& Documentation}{11}{15}
		\workpackage{Poster Preparation}{13}{15}
	\end{semestergantt}
	\caption{Final semester plan. All eight work packages are complete; the
	last two weeks were spent on the report and the poster.}
	\label{fig:gantt}
\end{figure}

\subsection{Description of Work Packages}
\begin{itemize}
	\item \textbf{Literature Review (W1--W3).} Survey of MEAM, EAM, and
	      neural-network potential work. Complete.
	\item \textbf{DFT Reference Calculations (W2--W5).} Quantum Espresso build,
	      pseudopotential management, and the EOS, elastic-constant, and
	      binary-pair runs. Complete.
	\item \textbf{MEAM Initialization Pipeline (W4--W7).} Auto-discovery and
	      merging of the library files, the DFT-to-MEAM analytic mapping, and
	      the unified twelve-element library. Complete.
	\item \textbf{NN Surrogate Training (W6--W9).} Parameter-space sampling,
	      parallel LAMMPS evaluation, and training of both the composition and
	      parameter surrogates. Complete.
	\item \textbf{Inverse Optimization (W8--W13).} Adam-based back-propagation
	      through the surrogate to fit the elastic targets, with the stability
	      limitations reported in Section~\ref{sec:nnip}. Complete.
	\item \textbf{Validation \& Testing (W10--W14).} Comparison against named
	      industrial alloys and the ASM/MatWeb external set. Complete.
	\item \textbf{Writing \& Documentation (W11--W15).} Final report and code
	      documentation.
	\item \textbf{Poster Preparation (W13--W15).} Poster for the
	      end-of-semester presentation.
\end{itemize}

\subsection{Code Availability}

The full pipeline, configuration generation, the LAMMPS and Quantum Espresso
drivers, and the neural-network surrogate is openly
available at \url{https://github.com/canndurmaz/PHYS400}.
